\section*{Chapter \ref{ch:b3:milky-way-galaxies} --- The Milky Way and Galaxies}

\begin{solution}{exo:b3:milky-way-galaxies:1}
(a)~$2\pi R_0 = 2\pi\times\qty{2.523e17}{km} = \qty{1.585e18}{km}$; at
\qty{229}{km.s^{-1}} this takes \qty{6.92e15}{s}, that is
\qty{219}{Myr}. (b)~$4540/219 = 20.7$: about twenty turns.
\end{solution}

\begin{solution}{exo:b3:milky-way-galaxies:2}
(a)~$\qty{1}{au}/\qty{1}{yr} = \num{1.496e8}/\num{3.156e7} =
\qty{4.74}{km.s^{-1}}$. (b)~$v = 4.74\times6.411\times8.178 =
\qty{248.5}{km.s^{-1}}$, within the review's
\qty[separate-uncertainty=true]{248 \pm 3}{km.s^{-1}}. (c)~$248.5 - 11 = \qty{237.5}{km.s^{-1}}$,
some \qty{8}{km.s^{-1}} above the \qty{229}{km.s^{-1}} of the red giants:
the two rest on different values of $R_0$ (8.178 and \qty{8.122}{kpc}),
and the LSR itself may stream by up to about \qty{15}{km.s^{-1}} relative
to a truly circular orbit.
\end{solution}

\begin{solution}{exo:b3:milky-way-galaxies:3}
(a)~$n = 10(1 - 0.6) = 4$: E4. (b)~$b/a = 1 - 7/10 = 0.3$. (c)~The class
uses the apparent axis ratio: a flattened ellipsoid seen along its short
axis has a circular outline and is classified E0.
\end{solution}

\begin{solution}{exo:b3:milky-way-galaxies:4}
$d = 10^{\mu/5 + 1}$~pc: the LMC \qty{49.6}{kpc}, Andromeda
\qty{783}{kpc}, NGC~4258 \qty{7.58}{Mpc}. A star of the Cloud has a
parallax of $1/\qty{49.6}{kpc} = \qty{0.020}{mas}$, the size of the best
catalogue's uncertainty for bright stars: a single parallax gives no
distance at all, and the Cloud is placed by other means (eclipsing
binaries), from which the Cepheid scale starts.
\end{solution}

\begin{solution}{exo:b3:milky-way-galaxies:5}
(a)~$\omega_0 = A - B = \qty{27.2}{km.s^{-1}.kpc^{-1}}$ and $V_0 =
27.2\times8.178 = \qty{222}{km.s^{-1}}$. (b)~$V'_0 = -(A + B) =
\qty{-3.4}{km.s^{-1}.kpc^{-1}}$: a gently falling curve. (c)~Both are
below the red-giant values (229 and \qty{-1.7}{km.s^{-1}} per
kiloparsec); the authors of that curve found the same, a $V_0$ about two
standard deviations lower from $R_0$ and these \omterm{def:b3:milky-way-galaxies:rotation-curve}{Oort constants}. Oort's
formulas assume circular orbits in an axisymmetric disc; non-zero $C$ and
$K$ show that the local velocity field is perturbed by the bar and the
spiral arms, and the local $A$ and $B$ carry those perturbations.
\end{solution}

\begin{solution}{exo:b3:milky-way-galaxies:6}
(a)~$\rho = V^2/(4\pi GR^2) = 229^2/(4\pi\times\num{4.301e-6}\times
8.178^2) = \num{1.45e7}$ solar masses per cubic kiloparsec, or
\num{0.0145} per cubic parsec. (b)~$M_\odot c^2 = \num{1.9884e30}\times
(\num{2.998e8})^2 = \qty{1.787e47}{J} = \qty{1.115e57}{GeV}$, and
$\qty{1}{pc^3} = (\qty{3.0857e18}{cm})^3 = \qty{2.938e55}{cm^3}$: the
ratio is \qty{37.96}{GeV.cm^{-3}}. Then $0.30/37.96 = \num{0.0079}$ solar
masses per cubic parsec. (c)~The flat-curve density counts all the mass
that holds the orbits, stars and gas included, as if it were spread in a
sphere; the dark matter is only part of it, about half here.
\end{solution}

\begin{solution}{exo:b3:milky-way-galaxies:7}
(a)~$W_I(\mathrm{LMC}) = -3.314\times1.4771 + 15.888 = \qty{10.993}{mag}$,
so $M_W = 10.993 - 18.477 = \qty{-7.484}{mag}$. (b)~$\mu = 25.00 + 7.484 =
32.484$ and $d = 10^{32.484/5 + 1}~\mathrm{pc} = \qty{31.4}{Mpc}$.
(c)~$\delta d/d = (\ln 10/5)\,\delta\mu = 0.4605\times0.077 =
\qty{3.5}{\percent}$ for one star, $3.5/\sqrt{50} = \qty{0.5}{\percent}$
for fifty, if their errors are independent (the error of the Cloud's own
distance adds to all of them alike).
\end{solution}

\begin{solution}{exo:b3:milky-way-galaxies:8}
$\mu = I - A_I - M_I = 26.00 - 0.05 + 4.05 = 30.00$, so $d =
10^{7}~\mathrm{pc} = \qty{10.0}{Mpc}$. An error of \qty{0.04}{mag} is
$0.4605\times0.04 = \qty{1.8}{\percent}$ in distance.
\end{solution}

\begin{solution}{exo:b3:milky-way-galaxies:9}
(a)~$d = cz/H_0$ gives $5\log_{10}d = 5\log_{10}(cz) - 5\log_{10}H_0$,
so $m_B - M_B - 25 = 5\log_{10}(cz) - 5\log_{10}H_0$; with
$\log_{10}(cz) = a_B + 0.2\,m_B$, $5\log_{10}H_0 = 5a_B + M_B + 25$, which
is the result. (b)~$\log_{10}H_0 = -3.8506 + 0.714158 + 5 = 1.8636$, so
$H_0 = \qty{73.04}{km.s^{-1}.Mpc^{-1}}$. (c)~$M_B = 5(\log_{10}67.4 -
5.714158) = \qty{-19.427}{mag}$: supernovae brighter by \qty{0.174}{mag},
hence all distances larger by a factor $10^{0.2\times0.174} = 1.084$, the
ratio $73.04/67.4$. The tension is a disagreement of eight percent in the
scale of distances.
\end{solution}

\begin{solution}{exo:b3:milky-way-galaxies:10}
(a)~A star at the edge, in a circular orbit, has $V^2/R \approx GM/R^2$
(spherical approximation). (b)~$R = \sqrt{M/(\pi\Sigma)}$, so $V^2 =
GM/R = G\sqrt{\pi\Sigma}\,M^{1/2}$ and $M = V^4/(\pi\Sigma G^2) \propto
V^4$; with a common mass-to-light ratio, $L \propto V^4$ too.
(c)~$(100/229)^4 = \num{0.036}$: the faster galaxy is 27 times more
massive. (d)~Same luminosity: the flux falls as $1/d^2$, so the fainter
one is twice as far.
\end{solution}

\begin{solution}{exo:b3:milky-way-galaxies:11}
(a)~$\qty{1}{km.s^{-1}} = \qty{1.0227}{kpc.Gyr^{-1}}$, so
\qty{229}{km.s^{-1}} is \qty{234.2}{kpc.Gyr^{-1}}. (b)~$2\pi R/V$:
\qty{107}{Myr} at \qty{4}{kpc}, \qty{322}{Myr} at \qty{12}{kpc}.
(c)~$\Delta\phi = Vt\,(1/4 - 1/12) = 234.2\times0.1667 = \qty{39.0}{rad}$,
6.2 turns: the arm would be a coil of six tight turns between 4 and
\qty{12}{kpc}, far tighter than observed arms, and tighter still at each
further rotation. (d)~In a density wave the pattern turns at one angular
speed, as a solid body; the stars and gas pass through it, so the shape
of the arms does not wind up.
\end{solution}

\begin{solution}{exo:b3:milky-way-galaxies:12}
(a)~$a = (4.2\times10^6\times16^2)^{1/3} = \qty{1024}{au}$; at
\qty{8178}{pc} it subtends $1024/8178 = \ang{;;0.125}$ (an au at one parsec
subtends one arcsecond). (b)~$R_S = 2GM/c^2 = 2\times\num{6.674e-11}
\times\num{8.35e36}/(\num{2.998e8})^2 = \qty{1.24e10}{m} = \qty{0.083}{au}$;
the diameter $2R_S$ subtends $0.166/8178 = \num{2.0e-5}$ arcsecond, about
\qty{20}{\micro as}. (c)~$R_S = \qty{4.02e-7}{pc}$, and $M/(\tfrac43\pi
R_S^3) = \num{1.5e25}$ solar masses per cubic parsec, forty times the
lower bound set by the motion and size of Sgr~A*: the bound comes within
a factor forty of the density of the black hole itself.
\end{solution}

\begin{solution}{pb:b3:milky-way-galaxies:1}

\textbf{Part I.}

\textbf{1.} $v_r = A\,d\sin2\ell$: $+A d$ at \ang{45}, 0 at \ang{90}, $-A d$ at
\ang{135}, 0 at \ang{180}.

\textbf{2.} $A = (15.4 + 15.2)/(2\times\qty{1.00}{kpc}) =
\qty{15.3}{km.s^{-1}.kpc^{-1}}$.

\textbf{3.} At \ang{90}, $v_t = d(B - A) = 4.74\times(-5.91)\times1.00 =
\qty{-28.02}{km.s^{-1}}$, so $B = 15.3 - 28.02 =
\qty{-12.72}{km.s^{-1}.kpc^{-1}}$.

\textbf{4.} At \ang{0}, $v_t = d(A + B) = 4.74\times0.55 =
\qty{2.61}{km.s^{-1}}$, so $B = 2.61 - 15.3 = \num{-12.69}$: the two
agree within \num{0.03}.

\textbf{5.} $\omega_0 = A - B = \qty{28.0}{km.s^{-1}.kpc^{-1}}$ and
$V_0 = 28.0\times8.18 = \qty{229}{km.s^{-1}}$.

\textbf{6.} $V'_0 = -(A + B) = \qty{-2.6}{km.s^{-1}.kpc^{-1}}$: gently falling.

\textbf{7.} $P = 2\pi/\omega_0$: $2\pi\times\qty{8.18}{kpc}/\qty{229.2}{km.s^{-1}}
= \qty{219}{Myr}$; $4540/219 = 20.7$ turns.

\textbf{Part II.}

\textbf{8.} The tangent point needs a line of sight that passes closer to the
centre than the Sun is, $\sin\ell > 0$ with $\ell < \ang{90}$ (or the
mirror quadrant, $\ang{270} < \ell < \ang{360}$, where the largest
velocity is the most negative). Outside the solar circle the distance
to the centre grows all along the line of sight: no point is singled
out, and the velocity alone does not say where the gas is.

\textbf{9.} $R = 8.18\sin\ell$: \qtylist{5.26;6.27;7.08;7.69}{kpc}.

\textbf{10.} $V = v_\mathrm{max} + V_0\sin\ell$: $82.9 + 147.2 = 230.1$,
$55.6 + 175.4 = 231.0$, $32.4 + 198.3 = 230.7$, $14.9 + 215.2 =
\qty{230.1}{km.s^{-1}}$.

\textbf{11.} $\Delta\nu = \nu v/c = \num{1.420406e9}\times82.9/299792 =
\qty{393}{kHz}$, towards lower frequencies: the gas recedes.

\textbf{12.} Within about \qty{3}{km.s^{-1}} of the red giants: between 5 and
\qty{8}{kpc} the curve is flat.

\textbf{13.} A point mass gives $V = V_0\sqrt{R_0/R} = 229\times\sqrt{8.18/5.26} =
\qty{286}{km.s^{-1}}$ at the first tangent point, much faster than the
\qty{230}{km.s^{-1}} found: the mass is not concentrated at the centre,
much of it lies between 5 and \qty{8}{kpc} and beyond.

\textbf{Part III.}

\textbf{14.} $\num{1.067e10} + 2\times\num{3.944e10} = \num{8.96e10}$ solar
masses, nearly twice the review's \num[separate-uncertainty=true]{5 \pm 1 e10}; the model's stars
may be too heavy, which would make the dark matter found below an
underestimate.

\textbf{15.} $V_b^2 = GMR^2/(R^2 + b^2)^{3/2} = \num{4.301e-6}\times\num{1.067e10}
\times400/(400.09)^{3/2}$, $V_b = \qty{47.9}{km.s^{-1}}$.

\textbf{16.} Thin: $(a + b)^2 = 30.80$, $V^2 = \num{4.301e-6}\times\num{3.944e10}
\times400/(430.8)^{3/2}$, $V = \qty{87.1}{km.s^{-1}}$. Thick:
$(a + b)^2 = 11.56$, $V = \qty{90.1}{km.s^{-1}}$.

\textbf{17.} $V_\mathrm{vis} = \sqrt{47.9^2 + 87.1^2 + 90.1^2} =
\qty{134.2}{km.s^{-1}}$.

\textbf{18.} At \qty{8.18}{kpc}: 74.8, 108.4 and \qty{127.8}{km.s^{-1}},
together \qty{183.5}{km.s^{-1}}, against $V_0 = \qty{229}{km.s^{-1}}$:
the visible mass falls short already at the Sun.

\textbf{19.} $\sqrt{G\times\num{8.96e10}/20} = \qty{138.8}{km.s^{-1}}$. The
model's \qty{134.2}{km.s^{-1}} is close, because at \qty{20}{kpc},
four to six disc scale lengths out, nearly all the visible mass lies
inside the orbit: the visible curve has become Keplerian.

\textbf{Part IV.}

\textbf{20.} $M = 20\times200^2/\num{4.301e-6} = \num{1.860e11}$ solar masses.

\textbf{21.} $M_\mathrm{vis} = 20\times134.2^2/\num{4.301e-6} = \num{8.37e10}$
solar masses.

\textbf{22.} $M_\mathrm{dark} = \num{1.860e11} - \num{8.37e10} = \num{1.02e11}$;
since both masses are $RV^2/G$ at the same $R$, the fraction is
$1 - (V_\mathrm{vis}/V)^2 = 1 - (134.2/200)^2 = \num{0.550}$.

\textbf{23.} $x = 20/14.8 = 1.351$, $\ln(2.351) - 1.351/2.351 = 0.8549 - 0.5748
= 0.2801$, $M = 4\pi\times\num{1.06e7}\times14.8^3\times0.2801 =
\num{1.21e11}$ solar masses; the dark fraction is
$1.21/(1.21 + 0.837) = \num{0.59}$. The halo is fitted to the whole
curve, which at \qty{20}{kpc} it overshoots: the measurements dip there
by about \qty{10}{km.s^{-1}} below the smooth model (\qty{209}{km.s^{-1}}).

\textbf{24.} $x = 8.18/14.8 = 0.553$, $\rho = \num{1.06e7}/(0.553\times1.553^2)
= \num{7.95e6}$ solar masses per cubic kiloparsec, \num{0.0080} per
cubic parsec, or \qty{0.30}{GeV.cm^{-3}}, the value the authors of the
fit quote.

\textbf{25.} $1 - 8.37/18.60 = \num{0.550}$: the dark-matter fraction of the
mass within \qty{20}{kpc} is
\textbf{\qty[text-series-to-math=true]{55}{\percent}}.
\end{solution}
